\minitopic{研读实验室：确证理论的坐标系}\index{回溯问题}\index{基于关系}\index{内在主义}\index{外在主义} 基础主义、融贯主义、内在主义、外在主义经常被排成四个互斥选项，这是分类错误。前两者回答支持结构：理由链如何终止或相互支撑；后两者回答决定因素：使信念得到确证的事实是否必须以某种方式对主体可及。一个理论可以是外在主义基础论，也可以是内在主义融贯论。先画坐标，再比较方案，能避免只记标签。 还要区分确证的对象。命题确证是主体的总证据支持$p$；信念确证是主体实际因为适当根据而相信$p$。一位学生有充分数据证明答案，却只因梦到选项C而选C：命题得到支持，信念未适当基于支持。反之，他可能以正常视觉相信屏幕为蓝，即使不能陈述视觉可靠率，外在主义会认为信念得到确证。 “有理由”也可指主体心理中拥有一项内容、客观上存在可取得证据、或规范上应该去寻找证据。数据库中有一份反驳报告但主体无权访问，它是否击败信念，取决于理论讨论个人确证、机构知识还是探究责任。教材案例必须标明信息范围。

\minitopic{回溯问题的三类答案}若每个有确证信念都要由另一个有确证信念支持，就出现回溯：链条终止、循环或无限延伸。基础主义允许某些基础信念以非推论方式获得初步确证；融贯主义把支持放在系统整体的解释与相互约束；无限主义认为理由链可以无重复地继续，而主体不必实际同时思考全部理由。 基础信念不是“没有原因的武断信念”。知觉经验、记忆呈现或内省可提供非推论根据，确证仍可被击败。温和基础主义允许基础支持是可错、有程度且需响应反证的。问题转为：经验如何具有理由地位，主体是否必须概念化经验，以及不同经验冲突时如何比较。 融贯也不只是逻辑一致。许多互不相干的假故事可以一致；强融贯要求解释联系、信息整合、反事实稳定和经验输入。若系统只靠内部互相支持，孤立于世界的完美小说仍可能融贯。融贯主义因而需要说明感觉与证言如何给系统施加非任意约束。 无限主义提醒我们，任何有限终点都可能被继续追问。但“总能再给一个理由”与“主体现在拥有无限理由”不同。若理由只在被问时临时构造，确证是否早已存在？无限链也可能整体偏离世界。其贡献在于挑战回溯只能在基础与循环间选择，而不是轻易提供完整心理模型。

\begin{argumentbox}[回溯论证]
（1）若信念B得到推论确证，就有理由R支持B；（2）R若要提供确证，本身也需适当认识地位；（3）这一要求可重复；所以理由结构要么终止、循环、无限延伸，要么产生怀疑结论。各理论并不否认重复，而是在说明哪一种结构能够真正传递或生成支持。
\end{argumentbox}

\minitopic{基于关系：理由如何真正发挥作用}主体拥有理由R且相信$p$，仍不保证他基于R相信$p$。基于关系至少包含反事实与因果维度：若R改变，主体的信念是否相应改变；R是否在形成或维持信念中发挥解释作用。单一测试都不充分，因为信念可能由多个理由过度决定，也可能在反思后改变基础。 设医生既看到可靠检测，又迷信当天星象；两者都支持她相信病人患病。若检测单独足够且主导判断，迷信理由未必污染信念；若她无论检测如何都随星象判断，可靠证据只是心理装饰。分析应问实际推理、注意分配和更新模式，而非只列主体口头能说出的理由。 主体也会事后编造理由。先由偏见形成信念，再寻找看似中立统计支持，形成动机性合理化。新证据可能后来真正成为维持基础，但不能回溯使早期信念得到确证。时间切片和因果历史都重要。 基于关系不必要求有意识推理。熟练驾驶者看见红灯便刹车，经验可直接调节信念与行动；儿童以照护者证言学习也未列出可靠性前提。把一切基于都要求主体陈述三段论，会不当排除大量知识。

\minitopic{击败者的分类与处理}反驳型击败者支持非$p$，例如另一份可靠测量显示温度不高；削弱型击败者不支持非$p$，却破坏原证据联系，例如知道温度计近期频繁失灵。两者都可降低确证，但响应不同：反驳型需要比较相反证据，削弱型需要重建来源可靠性。 心理击败者是主体实际接受并据此降低信心的信息；规范击败者可能主体未接受，却按其可得证据本应响应。一个学生看到标准答案与自己不同却无理由忽视，原信念受规范击败；一份锁在他无从得知的档案是否击败个人确证，则有争议。机构层面若有人负责接收该档案，答案可能不同。 误导性击败者虽然为假，却可合理降低信心。若实验室发布错误召回通知，研究者暂停依赖设备是理性的；后来确认通知发错，击败者被击败。理论若只允许真信息击败，就难以解释主体为何应响应可信但错误的警报。 高阶证据直接评价主体的能力或证据处理。你完成计算后得知自己服用了显著损害运算的药物，即使找不到具体哪一步错，也应降低确信。高阶证据是否总压过一阶证据没有统一答案：清楚逐步证明可能抵抗含混能力警告，反复独立发现错误则应显著降级。

\minitopic{内在与外在：新邪恶恶魔的压力}设受骗主体的全部经验和推理与我们相同，却因恶魔操纵几乎所有外部信念为假。他似乎在主观上同样理性；纯粹实际可靠性理论若说他没有任何确证，就遗漏这种内部无责与证据合宜。案例压力针对确证，不直接证明他拥有知识，因为知识仍要求真。 可靠主义者可区分客观保证与主观确证，或以正常世界、设计功能而非实际环境统计评价过程。这样保存受骗者的某种正面地位，却需要说明为何评价环境不是任意挑选。若以“本来应在的环境”为准，演化、设计和社会训练分别可能给出不同标准。 内在主义的反面压力来自儿童、动物与专家直觉。主体可以可靠辨认，却无法接触自己的过程可靠性；若确证要求对所有决定因素反思可及，就会引起更高阶回溯。一种折衷是：一阶知识可依赖外在能力，反思确证要求主体对显著风险有适当视角。折衷不是自动胜利，仍要解释两层关系与术语使用。

\minitopic{比较视角：闻见、思虑与学习结构}《荀子》讨论认识时常把感官接触、心的辨别、名的公共使用与长期学习放在同一实践中\parencite{xunzi2014}。这给当代分类一种压力：理由既不只是主体瞬时可见的命题，也不只是脱离训练的外部可靠过程。感官提供材料，辨别组织材料，师法与语言规范使判断可校正。 但把“心”直接译成当代内在主义，或把“师法”直接等同外在可靠机制，都过度。前者会忽略心在行动、情感与分类中的广泛作用，后者会忽略规范学习和公共辨正。更合适的比较问题是：一种认识实践如何同时安置信息输入、主体反思和社会校准？ 当代基础主义常从个体时刻的经验开始，荀子式学习图景则提醒我们，主体能把经验当作理由已经依赖概念、语言与训练。反过来，当代理论对击败、可靠与基于关系的细分，也能让“学习”不再是笼统德目：哪些反馈改变信念，何种权威可被纠正，错误如何暴露，都可具体分析。

\minitopic{综合案例：档案研究中的来源冲突}历史学家发现一封署名部长的信，纸张、墨水与措辞均符合年代，内容支持一项政策解释。她相信部长确曾下令。后来同事指出，档案馆同一批次曾混入高质量仿件。此信息不是直接证明信为假，而是削弱来源链。 命题确证方面，现有材料仍部分支持下令；信念确证方面，要看研究者是否因文件特征而相信，还是早已有政治立场并把信件当装饰。基础主义会问视觉与材料检测提供何种初步支持；融贯主义会比较其他会议记录、财政数据与政策结果；可靠主义审查鉴定流程在该档案环境的表现；内在主义追问研究者可得理由。 若鉴定员用独立化学检测确认墨水年代，却不知道伪造者使用旧墨，信念可能仍受环境运气影响。若另一封独立来信提到同一命令，融贯增强，但两封信若源自同一伪造工作室并不独立。支持数量必须结合依赖结构。 研究者最终应把结论写成与证据相称的形式：“现有书信与政策结果共同支持部长下令，但混档记录使作者身份尚待独立验证。”这不是逃避判断，而是分别报告一阶证据、击败信息与剩余不确定性。确证理论的实践价值正在于帮助我们精确降级，而非只给出相信或不信的二元命令。

\begin{tcolorbox}[breakable,colback=white,colframe=blue!30!black,title={本章研读任务},fonttitle=\bfseries]
为一个争议信念绘制支持图：标出非推论输入、推论理由、相互支持、外部可靠机制、主体可及内容和击败者。再说明信念产生基础与当前维持基础是否相同，并分别给出基础主义、融贯主义、内在主义和外在主义评价。
\end{tcolorbox}
